% =========================
% 4_implementation.tex --- what we built, on what, and what it costs
% =========================
\section{Implementation}
\label{sec:implementation}

We implement HALO in PyTorch as a single repository covering conversion, pretraining, memory
construction, and the evidence rule. The system is deliberately small where it matters: the trunk
holds 7.70\,M parameters and the only component fitted after Phase A is a rank-2 bilinear gate with
roughly 1.5\,k parameters.

\subsection{Codebase}
\label{sec:codebase}

Table~\ref{tab:implementation} lists the components. We write one converter per dataset --- 35 in
total --- rather than a generic ingest path, because every public corpus differs in column order,
unit convention, gravity handling, and subject encoding, and a generic reader hides exactly the
errors that later corrupt a physical feature. Each converter emits whole-recording sessions; the
grid builder owns all windowing, so no converter can silently choose its own window policy.

\begin{table}[t]
  \centering
  \caption{Implementation components. Line counts exclude generated artifacts.}
  \label{tab:implementation}
  \scriptsize
  \setlength{\tabcolsep}{3pt}
  \begin{tabular}{@{}lrp{0.44\columnwidth}@{}}
    \toprule
    Component & Lines & Contents \\
    \midrule
    Dataset converters & 11.8\,k & 35 per-dataset readers, unit and gravity normalization \\
    Corpus pipeline    &  6.9\,k & curation policy, windowing, three grid alignments, quality screens \\
    Model              &  3.9\,k & filterbank, sensor folding, conditioning, dual-branch trunk \\
    Phase-A training   &  7.6\,k & sampler, augmentation, masking, objectives, telemetry \\
    Evidence stage     & 14.3\,k & bank build, resolvability, gate, prediction rule, protocol \\
    Evaluation harness &  2.1\,k & baseline adapters, scoring, protocol fingerprinting \\
    Tests              &  9.3\,k & 578 tests across 45 files \\
    \bottomrule
  \end{tabular}
\end{table}

\subsection{Training and Memory Cost}
\label{sec:cost}

We train on one NVIDIA RTX 4090 (24\,GB). The reported Phase-A run completes 30,000 steps at batch
256 in 33 minutes of wall clock and peaks at 4.36\,GiB of allocated device memory, so a full
pretraining run fits comfortably on a single consumer GPU. Mixed-precision training is enabled and
we compile only the transformer; compiling the encoder specializes on each batch's changing
unique-text cardinality and runs slower.

Memory construction is the expensive artifact. The schema-4 bank stores one row per patch and
sensor, which is 7.65\,M rows for a 250,000-window archive and 6.9\,GB on disk in half precision.
Two engineering choices keep the evidence path tractable. First, we do not store a 384-dimensional
sensor descriptor on every row: roughly $10^5$ rows are drawn from about 40 distinct sensor
descriptions, so we store a de-duplicated table and an index, evaluate the gate once per distinct
description, and gather. This is exact, not approximate, because the gate reads only text. Second,
we cache stream metadata and text encodings per configuration, so sentence-embedding inference never
runs inside a query loop.

\subsection{Artifact Binding}
\label{sec:guards}

A retrieval system built on a cached bank fails in a specific way: an encoder is retrained, the bank
is not rebuilt, and every downstream number silently describes a stale embedding space. We make that
failure impossible to reach rather than easy to notice.

Every artifact carries a fingerprint of what produced it, and every consumer checks it before doing
any work. The bank records the encoder checkpoint hash, the label vocabulary, and behavioral probes
of both the embedding path and the sentence-embedding path; loading it re-runs those probes against
the live encoder and refuses a mismatch. The gate records the resolvability table hash, the
checkpoint that produced it, and the bank it is bound to. The evaluation protocol is itself
fingerprinted over the source files that define scoring behavior, so a change to the protocol
invalidates cached results instead of blending two protocols into one table. The quality screens are
loaded with the same discipline: a missing or stale exclusion cache is a hard failure, because
silently re-admitting byte-identical stale-buffer windows would corrupt both training and scoring.

These guards are load-bearing rather than decorative. During development they rejected a legacy
resolvability table that pooled evaluation datasets into a training-only measurement, and they
currently refuse the gate artifact described in Section~\ref{sec:label_efficiency} whenever it is
paired with a bank from a different checkpoint.

\subsection{Reproducibility}
\label{sec:reproducibility}

Every quantitative statement in this paper is bound to a named artifact through a checked-in evidence
ledger that records the producing command, the source-data manifest, the seed, the checkpoint
fingerprint, and the raw output. Values without that chain are not printed as results, which is why
several quantities from an earlier draft of this work do not appear here. We will release the
converters, the corpus pipeline, the model, the evidence stage, and the evaluation harness.
